\subsection{Multiplicities and secant sequences}

PROPOSITION 7. --- Assume A is local. Let s be an integer ≥ 1 and, for \(1 \leqslant i \leqslant s\) , let \(\delta_{i}\) be an integer \textgreater{} 0, \(x_{i}\) an element of \(m_{A}^{\delta_{i}}\) , and \(\xi_{i}\) its class in \(m_{A}^{\delta_{i}}/m_{A}^{\delta_{i}+1}\) . Assume that \((x_{1}, \ldots, x_{s})\) is a secant sequence for A. Let x denote the ideal of A generated by \((x_{1}, \ldots, x_{s})\) . Then one has \(e(A/x) \geqslant \delta_{1} \ldots \delta_{s} \cdot e(A)\) with equality if \((\xi_{1}, \ldots, \xi_{s})\) is a completely secant sequence for gr(A).

Set B = A/x, and consider the formal series

\[
\mathrm{H} _ {\mathrm{A}} = \sum_ {n \geqslant 0} \operatorname{long} (\mathfrak {m} _ {\mathrm{A}} ^ {n} / \mathfrak {m} _ {\mathrm{A}} ^ {n + 1}). \mathrm{T} ^ {n}, \quad \mathrm{H} _ {\mathrm{B}} = \sum_ {n \geqslant 0} \operatorname{long} (\mathfrak {m} _ {\mathrm{B}} ^ {n} / \mathfrak {m} _ {\mathrm{B}} ^ {n + 1}). \mathrm{T} ^ {n}
\]

and \(\mathrm{H}_{\mathrm{B}}^{(s)} = (1 - \mathrm{T})^{-s}\mathrm{H}_{\mathrm{B}}\) . According to prop. 6 of § 4, no. 5, one has in \(\mathbf{Z}[[T]]\) , the inequality

\[
\mathrm{H} _ {\mathrm{B}} ^ {(s)} \geqslant \left(\prod_ {i = 1} ^ {s} \frac {1 - \mathrm{T} ^ {\delta_ {i}}}{1 - \mathrm{T}}\right) \mathrm{H} _ {\mathrm{A}},
\]

and there is equality when the sequence \((\xi_{1}, \ldots, \xi_{s})\) is completely secant. But

\[
\mathrm{R(T)} = \prod_ {i = 1} ^ {s} \frac {1 - \mathrm{T} ^ {\delta_ {i}}}{1 - \mathrm{T}}
\]

is a polynomial in \(\mathbf{Z}[\mathrm{T}]\) such that \(\mathbf{R}(1) = \delta_1 \ldots \delta_s\) . Put \(\dim(\mathbf{A}) = d\) ; one has \(\dim(\mathbf{B}) = d - s\). By th. 2 of §4, no. 3, there exist elements \(\mathbf{R}_{\mathbf{A}}\) and \(\mathbf{R}_{\mathbf{B}}\) of \(\mathbf{Z}[\mathrm{T}, \mathrm{T}^{-1}]\) such that

\[
\mathrm{H} _ {\mathrm{A}} = (1 - \mathrm{T}) ^ {- d} \mathrm{R} _ {\mathrm{A}} (\mathrm{T}), \quad \mathrm{R} _ {\mathrm{A}} (1) = e (\mathrm{A}),
\]

\[
\mathrm{H} _ {\mathrm{B}} = (1 - \mathrm{T}) ^ {- d + s} \mathrm{R} _ {\mathrm{B}} (\mathrm{T}), \quad \mathrm{R} _ {\mathrm{B}} (1) = e (\mathrm{A} / x).
\]

Thus we have

\[
(1 - \mathrm{T}) ^ {- d} \mathrm{R} _ {\mathrm{B}} (\mathrm{T}) = \mathrm{H} _ {\mathrm{B}} ^ {(s)} \geqslant \left(\prod_ {i = 1} ^ {s} \frac {1 - \mathrm{T} ^ {\delta_ {i}}}{1 - \mathrm{T}}\right) \mathrm{H} _ {\mathrm{A}} = (1 - \mathrm{T}) ^ {- d} \mathrm{R} (\mathrm{T}) \mathrm{R} _ {\mathrm{A}} (\mathrm{T}),
\]

and equality holds if the sequence \((\xi_{1},\ldots,\xi_{s})\) is completely secant. We conclude by lemma 2 of § 4, no. 1.

Remark. --- Conversely, one can show (cf. p. 103, exerc. 4) that if A is regular and \(e(\mathrm{A}/\mathfrak{x})=\delta_{1}\ldots\delta_{s}\), then the sequence \((\xi_{1},\ldots,\xi_{s})\) is completely secant.

Example. --- Let A be the formal power series ring \(k[[X_1, ..., X_n]]\) over a field \(k\); let \(F_1, ..., F_s\) be elements of A, let \(\mathfrak{a}\) be the ideal they generate, and set \(B = A/\mathfrak{a}\). Denote by \(P_1, ..., P_s \in k[X_1, ..., X_n]\) the initial forms of the series \(F_1, ..., F_s\) and \(\delta_1, ..., \delta_s\) their respective degrees. If the sequence \(F_1, ..., F_s\) is secant in A, we have \(e(B) \geqslant \delta_1 ... \delta_s\) ; if the sequence \(P_1, ..., P_s\) is completely secant in the ring \(k[X_1, ..., X_n]\) , one has \(e(B) = \delta_1 ... \delta_s\) .

Consider for example the ring \(\mathbf{B} = k[[\mathrm{X},\mathrm{Y}]] / \mathfrak{a}\) , where \(\mathfrak{a}\) is generated by \(\mathrm{X}^2 + \mathrm{Y}^3\) and \(\mathrm{X}^2 + \mathrm{Y}^4\) ; the previous inequality gives \(e(\mathbf{B})\geqslant 4\); noting that \(\mathfrak{a}\) is generated by the elements \(\mathrm{X}^2 + \mathrm{Y}^3\) and \(\mathrm{Y}^4 -\mathrm{Y}^3\) for which the sequence of initial forms is completely secant, one obtains \(e(\mathbf{B}) = 6\) .
% semantic-fragments: legacy-input-seam
\relax{}
